\chapter{Машинно-зависимые оптимизации}

Машинно-зависимые оптимизации используют особенности конкретного
процессора: его систему команд, ширину регистров, задержки отдельных
команд. После сравнения по 8 байт поиск совпадений
по-прежнему остаётся главной частью работы сжатия, и его можно ускорить
векторными командами, сравнивающими 32 или 64~байта за одну операцию.
Такие команды есть не у всех процессоров x86-64, поэтому программа выбирает
реализацию во время выполнения.

\section{Сравнение командами AVX-512BW}
\label{sec:avx512}

Расширение AVX-512BW имеет 512-битные
регистры и команды сравнения байт, результат которых "--- 64-битная маска.
Функция \texttt{match\_length\_avx512} (листинг~\ref{lst:avx512})
загружает по 64~байта кандидата и упреждающего буфера, получает командой
\texttt{vpcmpneqb} маску неравных байт, и номер первого единичного бита
маски (\texttt{tzcnt}) равен длине совпадения. Последний неполный блок
загружается с маской (\texttt{\_mm512\_maskz\_loadu\_epi8}): такие загрузки не
обращаются к байтам вне маски, поэтому граница \texttt{limit} никогда не
нарушается. Совпадение длиной 258~байт проверяется за 5~шагов вместо 33.

\begin{lstlisting}[caption={Сравнение командами AVX-512BW},label={lst:avx512}]
__attribute__ ((target ("avx512bw,avx512vl")))
static inline size_t
match_length_avx512 (const uint8_t *candidate, const uint8_t *cursor,
                     size_t limit)
{
  size_t len = 0;
  while (len + 64 <= limit)
    {
      __m512i candidate_bytes = _mm512_loadu_si512 (candidate + len);
      __m512i cursor_bytes = _mm512_loadu_si512 (cursor + len);
      __mmask64 unequal = _mm512_cmpneq_epi8_mask (candidate_bytes,
                                                   cursor_bytes);
      if (unequal != 0)
        return len + (size_t) __builtin_ctzll (unequal);
      len += 64;
    }
  if (len < limit)
    {
      __mmask64 tail = (1ULL << (limit - len)) - 1;
      __m512i candidate_bytes = _mm512_maskz_loadu_epi8 (tail, candidate + len);
      __m512i cursor_bytes = _mm512_maskz_loadu_epi8 (tail, cursor + len);
      __mmask64 unequal = _mm512_mask_cmpneq_epi8_mask (tail, candidate_bytes,
                                                        cursor_bytes);
      if (unequal != 0)
        return len + (size_t) __builtin_ctzll (unequal);
    }
  return limit;
}
\end{lstlisting}

\textbf{Результаты} (таблица~\ref{tab:res-avx}). На hppavillion "---
единственном компьютере с AVX-512 "--- сжатие корпуса ускорилось в
1,21~раза при 19~\% меньшем числе команд. Выигрыш меньше, чем у сравнения
по 8 байт: тот шаг уже убрал основную побайтовую работу, а 64-байтный шаг
помогает только на совпадениях длиннее 8~байт. Файл mixed.dat с длинными
повторами ускорился в 1,43~раза, nci "--- в 1,25~раза, текстовые файлы "---
примерно в 1,15~раза. На debian, где AVX-512 нет, программа выполняет
прежний код сравнения по 8 байт, и разница (1,02$\times$) находится в
пределах шума.

\section{Сравнение командами AVX2}
\label{sec:avx2}

Процессор Ryzen 7 5700U компьютера debian не поддерживает AVX-512, но
имеет AVX2 с 256-битными регистрами. Функция \texttt{match\_length\_avx2}
(листинг~\ref{lst:avx2}) сравнивает по 32~байта: команда
\texttt{vpcmpeqb} даёт по байту на каждую пару, \texttt{vpmovmskb} собирает
из них 32-битную маску равных байт, и первый нулевой бит маски,
\texttt{tzcnt(\textasciitilde mask)}, равен длине совпадения. Остаток короче
32~байт обрабатывает функция сравнения по 8 байт, так что за границу
\texttt{limit} программа не читает.

\begin{lstlisting}[caption={Сравнение командами AVX2},label={lst:avx2}]
__attribute__ ((target ("avx2,bmi")))
static inline size_t
match_length_avx2 (const uint8_t *candidate, const uint8_t *cursor,
                   size_t limit)
{
  size_t len = 0;
  while (len + 32 <= limit)
    {
      __m256i candidate_bytes =
        _mm256_loadu_si256 ((const __m256i *) (candidate + len));
      __m256i cursor_bytes =
        _mm256_loadu_si256 ((const __m256i *) (cursor + len));
      uint32_t equal = (uint32_t)
        _mm256_movemask_epi8 (_mm256_cmpeq_epi8 (candidate_bytes,
                                                 cursor_bytes));
      if (equal != UINT32_MAX)
        return len + (size_t) _tzcnt_u32 (~equal);
      len += 32;
    }
  return len + match_length (candidate + len, cursor + len, limit - len);
}
\end{lstlisting}

Команды BMI2 \texttt{pext} и \texttt{pdep} рассматривались для побитового
ввода-вывода, но были отвергнуты: на Zen~2 они выполняются микрокодом за
сотни тактов. Это пример того, что машинно-зависимая
оптимизация должна учитывать не только наличие команды, но и её стоимость на
конкретном процессоре.

\textbf{Результаты.} На debian сжатие корпуса ускорилось в 1,13~раза
(таблица~\ref{tab:res-avx}), но почти весь выигрыш даёт один файл,
mixed.dat (1,44$\times$), с длинными совпадениями. На файлах Silesia в
сумме выигрыш около 2~\%: nci, samba и xml ускорились на 6--13~\%, а dickens,
osdb, sao, x-ray и random.dat замедлились на 7--10~\%
(рисунок~\ref{fig:paths-debian}). Вариант, проверяющий первые 8 байт
одним 64-битным сравнением перед 32-байтным циклом, этих файлов не
улучшил и в сумме оказался медленнее (15,7~млрд тактов), поэтому в
программу не вошёл.

\begin{table}[H]
\caption{Результаты векторного сравнения, млрд тактов на сжатие всего корпуса}
\label{tab:res-avx}
\begin{tabletext}
\centering
\begin{tabular}{|l|r|r|r|r|}
\hline
\multirow{2}{*}{\textbf{Шаг}} & \multicolumn{2}{c|}{\textbf{debian (AVX2)}} & \multicolumn{2}{c|}{\textbf{hppavillion (AVX-512)}} \\
\cline{2-5}
 & \textbf{такты} & \textbf{ускорение} & \textbf{такты} & \textbf{ускорение} \\
\hline
04 сравнение по 8 байт & 16,65 & "--- & 17,99 & "--- \\
05 AVX-512BW & 16,29 & 1,02 (не используется) & 14,90 & \textbf{1,21} \\
06 AVX2 & 14,40 & \textbf{1,13} & 14,69 & 1,01 (не используется) \\
\hline
\end{tabular}
\end{tabletext}
\end{table}

\section{Выбор реализации во время выполнения}
\label{sec:dispatch}

Программа должна запускаться на любом процессоре x86-64, поэтому ключ
\texttt{-march=native} не используется. Вместо этого цикл сжатия
описан один раз как функция \texttt{compress\_stream} с атрибутом
\texttt{always\_inline} и компилируется трижды: обычная версия
\texttt{compress\_word} и версии \texttt{compress\_avx2} и
\texttt{compress\_avx512} с атрибутом \texttt{target}. В каждой
копии векторная функция сравнения встроена в цикл, так что на каждого
кандидата нет косвенного вызова. Функция \texttt{match\_kind}
(листинг~\ref{lst:dispatch}) один раз на файл проверяет возможности
процессора встроенной функцией \texttt{\_\_builtin\_cpu\_supports} и
выбирает самую широкую доступную реализацию. Переменная окружения
\texttt{MZ77\_MATCH=word} или \texttt{MZ77\_MATCH=avx2} принудительно
выбирает реализацию; это позволило измерить все варианты на одном
компьютере. На других архитектурах векторный код не компилируется вовсе.

\begin{lstlisting}[caption={Выбор реализации сравнения},label={lst:dispatch}]
static enum MatchKind
match_kind (void)
{
  const char *forced = getenv ("MZ77_MATCH");
  if (forced && strcmp (forced, "word") == 0)
    return MATCH_WORD;
  __builtin_cpu_init ();
  int has_avx2 = __builtin_cpu_supports ("avx2")
    && __builtin_cpu_supports ("bmi");
  if (forced && strcmp (forced, "avx2") == 0 && has_avx2)
    return MATCH_AVX2;
  if (__builtin_cpu_supports ("avx512bw")
      && __builtin_cpu_supports ("avx512vl"))
    return MATCH_AVX512;
  if (has_avx2)
    return MATCH_AVX2;
  return MATCH_WORD;
}
\end{lstlisting}

\textbf{Сложность.} Векторное сравнение уменьшает число шагов сравнения
кандидата до $\lceil M/w \rceil$ при $w = 32$ (AVX2) или $w = 64$
(AVX-512): оценка~\eqref{eq:after-15} сохраняет вид, но с большим $w$.

\textbf{Сравнение реализаций.} На рисунках~\ref{fig:paths-hpp}
и~\ref{fig:paths-debian} показано ускорение каждой векторной реализации
относительно сравнения по 8 байт на каждом файле корпуса
(\texttt{perf stat}, программа шага~06). Неожиданный результат: на
hppavillion реализация AVX2, выбранная принудительно, в сумме на 3~\%
быстрее реализации AVX-512, которую программа выбирает по умолчанию
(14,38 против 14,83~млрд тактов), и не хуже неё почти на всех файлах.
Причина в работе не установлена; можно лишь отметить, что большинство
кандидатов различаются уже в первых байтах, и на них ширина сравнения не
важна. Порядок выбора (AVX-512 первым) сохранён таким,
каким он был измерен; его изменение "--- возможное продолжение работы.

\begin{figure}[H]
\centering
\begin{tikzpicture}
\begin{axis}[
  ybar, width=\textwidth, height=7cm, bar width=5pt,
  ylabel={Ускорение относительно сравнения по 8 байт}, ymin=0.8, ymax=1.6,
  xtick=data, xticklabels from table={data/paths-hppavillion.dat}{file},
  xticklabel style={rotate=60, anchor=east, font=\small\ttfamily},
  legend style={at={(0.98,0.98)}, anchor=north east, font=\small},
  enlarge x limits=0.04, grid=major, grid style={gray!25},
]
\addplot[fill=gray!40] table[x=index, y expr=\thisrow{word}/\thisrow{avx512}] {data/paths-hppavillion.dat};
\addlegendentry{AVX-512BW}
\addplot[fill=gray!85] table[x=index, y expr=\thisrow{word}/\thisrow{avx2}] {data/paths-hppavillion.dat};
\addlegendentry{AVX2}
\end{axis}
\end{tikzpicture}
\caption{Ускорение векторных реализаций сравнения по файлам, hppavillion (такты)}
\label{fig:paths-hpp}
\end{figure}

\begin{figure}[H]
\centering
\begin{tikzpicture}
\begin{axis}[
  ybar, width=\textwidth, height=7cm, bar width=8pt,
  ylabel={Ускорение относительно сравнения по 8 байт}, ymin=0.8, ymax=1.6,
  xtick=data, xticklabels from table={data/paths-debian.dat}{file},
  xticklabel style={rotate=60, anchor=east, font=\small\ttfamily},
  legend style={at={(0.98,0.98)}, anchor=north east, font=\small},
  enlarge x limits=0.04, grid=major, grid style={gray!25},
]
\addplot[fill=gray!85] table[x=index, y expr=\thisrow{word}/\thisrow{avx2}] {data/paths-debian.dat};
\addlegendentry{AVX2}
\draw[dashed] (axis cs:-1,1) -- (axis cs:16,1);
\end{axis}
\end{tikzpicture}
\caption{Ускорение реализации AVX2 по файлам, debian (такты)}
\label{fig:paths-debian}
\end{figure}

\screenshot{match-paths}{Сжатие корпуса с разными реализациями сравнения на hppavillion}{fig:match-paths}{%
На hppavillion, сборка шага 06 (\texttt{bench/mz77-bench}), три команды
подряд:\\
\texttt{taskset -c 3 perf stat -e cycles:u bench/mz77-bench --mode=compress --repeat=1 bench/corpus/*/*}\\
затем то же с \texttt{MZ77\_MATCH=avx2} и \texttt{MZ77\_MATCH=word} в начале
строки. Снять три итоговых числа \texttt{cycles:u}.}
